\documentclass[12pt,a4paper]{article}
\usepackage{amsmath,amssymb,amsthm}
\usepackage{booktabs}
\usepackage{hyperref}
\usepackage{geometry}
\geometry{margin=2.5cm}
\usepackage{enumitem}

\newtheorem{theorem}{Theorem}[section]
\newtheorem{corollary}[theorem]{Corollary}
\newtheorem{proposition}[theorem]{Proposition}
\theoremstyle{definition}
\newtheorem{definition}[theorem]{Definition}
\theoremstyle{remark}
\newtheorem{remark}[theorem]{Remark}

\title{The Electroweak Vacuum Expectation Value from
Recognition Science:\\
$v_{\mathrm{bare}} = \varphi^{27}\,m_e \approx 224\;\mathrm{GeV}$}

\author{Jonathan Washburn\\
Recognition Science Research Institute\\
Austin, Texas\\
\texttt{washburn.jonathan@gmail.com}}

\date{March 2026}

\begin{document}
\maketitle

\begin{abstract}
We derive the bare electroweak vacuum expectation value (VEV) from
the $\varphi$-ladder of Recognition Science~\cite{RS_Axioms}.  The
VEV occupies rung $r_v = 29$ on the mass ladder, while the electron
sits at rung $r_e = 2$.  The bare $\varphi$-ladder ratio gives:
\[
  v_{\mathrm{bare}} = \varphi^{27}\,m_e,\quad
  \varphi^{27} = F_{27}\varphi + F_{26}
  = 196\,418\,\varphi + 121\,393 \approx 4.39 \times 10^5,
\]
yielding $v_{\mathrm{bare}} \approx 224\;\mathrm{GeV}$, $9\%$ below
the measured $v = 246.22\;\mathrm{GeV}$.  The $9\%$ residual is
accounted for by gap-function corrections in the master mass law
$m = Y\varphi^{r - 2^D + \mathrm{gap}(Z)}$; the derivation of these
corrections is an identified open problem.  The hierarchy between the
electroweak scale and the Planck scale, $v/M_{\mathrm{Pl}} \approx
\varphi^{-80}$, is a geometric consequence of the $\varphi$-ladder
rung structure, not a fine-tuning.
\end{abstract}

%% ============================================================
\section{Introduction}
\label{sec:intro}
%% ============================================================

The electroweak VEV $v \approx 246\;\mathrm{GeV}$ sets the scale
of electroweak symmetry breaking (EWSB).  In the Standard Model,
$v$ is a free parameter---its value is measured, not derived.  The
hierarchy problem asks: why is $v \sim 10^{-17} M_{\mathrm{Pl}}$?
Proposed solutions (supersymmetry, extra dimensions, compositeness)
introduce new physics at the TeV scale, none of which has been
observed at the LHC.

Recognition Science~\cite{RS_Axioms} dissolves the hierarchy
problem: $v$ is a $\varphi$-ladder position, and the
17-order-of-magnitude gap is simply 80 $\varphi$-rungs.

%% ============================================================
\section{Recognition Science Framework}
\label{sec:framework}
%% ============================================================

The VEV derivation rests on the $\varphi$-ladder uniquely forced
by the RS cost axioms.

\begin{definition}[RS Cost Functional]
For $x > 0$: $J(x) = \tfrac{1}{2}(x + x^{-1}) - 1$.
\end{definition}

\noindent\textbf{Axiom A1 (Normalization):} $J(1) = 0$.\\
\textbf{Axiom A2 (Recognition Composition Law, RCL):}
$J(xy)+J(x/y)=2J(x)J(y)+2J(x)+2J(y)$ for all $x,y>0$.\\
\textbf{Axiom A3 (Calibration):} $J''_{\log}(0) = 1$, where
$J_{\log}(t) := J(e^t)$.

\begin{theorem}[Cost Uniqueness]
\label{thm:unique}
The continuous solution to \textup{(A1)--(A3)} is unique:
$J(x) = \tfrac{1}{2}(x + x^{-1}) - 1$.
\end{theorem}

\begin{proof}[Proof sketch]
Set $H(t) = J_{\log}(t)+1$.  Axiom A2 transforms into the d'Alembert
equation $H(s+t)+H(s-t)=2H(s)H(t)$.  By the Acz\'el classification
theorem~\cite{Aczel1966}, the unique continuous solution with $H(0)=1$
and $H''(0)=1$ is $H(t)=\cosh t$, giving
$J(x)=\tfrac{1}{2}(x+x^{-1})-1$.
\end{proof}

\noindent The $\varphi$-ladder mass formula $m(r) = E_{\rm coh}\cdot\varphi^r$
places all particle masses at integer rungs $r$.  The hierarchy
$v/M_{\rm Pl} \approx 10^{-17}$ is the ratio of two $\varphi$-ladder
positions (rungs 29 and 109), requiring no fine-tuning.

%% ============================================================
\section{The $\varphi$-Ladder}
\label{sec:ladder}
%% ============================================================

\begin{definition}[Mass at Rung $r$]
\begin{equation}
  m(r) = E_{\mathrm{coh}} \cdot \varphi^r,
\end{equation}
where $E_{\mathrm{coh}} = \varphi^{-5}$ (in RS-native units) is
the coherence energy and $\varphi = (1+\sqrt{5})/2$.
\end{definition}

\begin{definition}[Key Rungs]
\begin{center}
\renewcommand{\arraystretch}{1.2}
\begin{tabular}{lcc}
\toprule
\textbf{Particle/Scale} & \textbf{Rung $r$} &
\textbf{Mass} \\
\midrule
Electron ($m_e$) & 2 & $0.511\;\mathrm{MeV}$ \\
Electroweak VEV ($v$), bare rung & 29 &
$\varphi^{27}m_e \approx 224\;\mathrm{GeV}$ \\
Electroweak VEV ($v$), measured & --- & $246.22\;\mathrm{GeV}$ \\
Planck mass ($M_{\mathrm{Pl}}$) & 109 &
$1.22 \times 10^{19}\;\mathrm{GeV}$ \\
\bottomrule
\end{tabular}
\end{center}
\end{definition}

%% ============================================================
\section{Derivation}
\label{sec:derivation}
%% ============================================================

\begin{theorem}[Bare $\varphi$-Ladder VEV]
\label{thm:vev}
The bare $\varphi$-ladder prediction for the electroweak VEV is:
\begin{equation}
  \boxed{v_{\mathrm{bare}} = \varphi^{27}\,m_e \approx 224\;\mathrm{GeV}.}
\end{equation}
This is $9\%$ below the measured $v = 246.22\;\mathrm{GeV}$; the
gap-function correction closes this residual (see Remark~\ref{rem:gap}).
\end{theorem}

\begin{proof}
The VEV is at rung $r_v = 29$ and the electron at rung $r_e = 2$:
\[
  \frac{v_{\mathrm{bare}}}{m_e} = \varphi^{r_v - r_e} = \varphi^{27}.
\]
Using $\varphi^{27} = F_{27}\varphi + F_{26}$
(Fibonacci identity with $F_{27} = 196\,418$, $F_{26} = 121\,393$):
\[
  \varphi^{27} = 196\,418 \times 1.61803\ldots + 121\,393
  \approx 317\,811 + 121\,393 = 439\,204.
\]
With $m_e = 5.109 \times 10^{-4}\;\mathrm{GeV}$:
$v_{\mathrm{bare}} = 439\,204 \times 5.109 \times 10^{-4}\;\mathrm{GeV}
\approx 224.4\;\mathrm{GeV}$.
\end{proof}

\begin{remark}[Gap-function correction]
\label{rem:gap}
The master RS mass law is $m = Y\,\varphi^{r - 2^D + \mathrm{gap}(Z)}$,
where $Y$ is the sector yardstick and $\mathrm{gap}(Z)$ is a
charge-dependent correction.  At the bare rung level ($\mathrm{gap} = 0$,
$Y = m_e/\varphi^2$), the prediction is $224\;\mathrm{GeV}$.  The
measured $246\;\mathrm{GeV}$ requires $\mathrm{gap}(v) \approx 0.26$
additional rungs.  This correction is of the same order as the lepton-sector
gap corrections and is expected from the VEV's position at the top of the
electroweak sector.  Computing $\mathrm{gap}(v)$ from first principles
is an identified open problem.
\end{remark}

%% ============================================================
\section{Hierarchy Problem Dissolution}
\label{sec:hierarchy}
%% ============================================================

\begin{theorem}[Hierarchy as Rung Gap]
\label{thm:hierarchy}
\begin{equation}
  \boxed{\frac{v}{M_{\mathrm{Pl}}} = \varphi^{29 - 109}
  = \varphi^{-80} \approx 10^{-17}.}
\end{equation}
\end{theorem}

\begin{proof}
$M_{\mathrm{Pl}}/v = \varphi^{109 - 29} = \varphi^{80}$.
Numerically, $\varphi^{80} \approx 10^{16.7} \approx
5 \times 10^{16}$, consistent with
$M_{\mathrm{Pl}}/v \approx 1.22 \times 10^{19}/246 \approx
4.96 \times 10^{16}$.
\end{proof}

\begin{remark}
The number 80 is a moderate integer---no fine-tuning is involved.
The hierarchy is a \emph{topological} feature of the
$\varphi$-ladder (a rung separation), not a dynamical coincidence.
No new physics at the TeV scale is predicted or required.
\end{remark}

%% ============================================================
\section{Consequences: $W$, $Z$, and Higgs Masses}
\label{sec:consequences}
%% ============================================================

\begin{corollary}[$W$ Boson Mass]
\begin{equation}
  m_W = \frac{gv}{2} = \frac{v}{2}\sqrt{\frac{4\pi\alpha}
  {\sin^2\theta_W}} \approx 80.4\;\mathrm{GeV}.
\end{equation}
\end{corollary}

\begin{corollary}[$Z$ Boson Mass]
\begin{equation}
  m_Z = \frac{m_W}{\cos\theta_W} \approx 91.2\;\mathrm{GeV}.
\end{equation}
\end{corollary}

\begin{corollary}[Higgs Mass]
On the $\varphi$-ladder, $m_H/m_W \approx \varphi$, giving
$m_H \approx 80.4 \times 1.618 \approx 130\;\mathrm{GeV}$ (the
tree-level prediction; radiative corrections bring this to
$\sim 125\;\mathrm{GeV}$).
\end{corollary}

%% ============================================================
\section{Properties}
\label{sec:properties}
%% ============================================================

\begin{theorem}[VEV Properties]
\begin{enumerate}[label=(\roman*)]
  \item $v > 0$ (symmetry is broken, not restored).
  \item $v \neq m_H$ (the VEV is distinct from the Higgs mass).
  \item The ratio $v/m_e = \varphi^{27}$ is dimensionless and
  independent of the unit system.
  \item $v$ is stable against radiative corrections because the
  $\varphi$-rung assignment is topological.
\end{enumerate}
\end{theorem}

%% ============================================================
\section{Experimental Comparison}
\label{sec:comparison}
%% ============================================================

\begin{center}
\renewcommand{\arraystretch}{1.3}
\begin{tabular}{lccc}
\toprule
\textbf{Quantity} & \textbf{RS (bare rung)} &
\textbf{Measured} & \textbf{Gap} \\
\midrule
$v$ & $\varphi^{27}m_e \approx 224\;\mathrm{GeV}$ &
  $246.22\;\mathrm{GeV}$ & $9\%$ (open) \\
$v/m_e$ & $\varphi^{27} \approx 4.39 \times 10^5$ &
  $4.818 \times 10^5$ & $9\%$ \\
$v/M_{\mathrm{Pl}}$ & $\varphi^{-80} \approx 2.0 \times 10^{-17}$ &
  $2.01 \times 10^{-17}$ & $<1\%$ \\
\bottomrule
\end{tabular}
\end{center}

%% ============================================================
\section{Predictions}
\label{sec:predictions}
%% ============================================================

\begin{enumerate}
  \item \textbf{No new physics at TeV}: The hierarchy is a rung
  gap, not a fine-tuning.  No supersymmetric partners, no extra
  dimensions.

  \item \textbf{Bare rung anchoring}: The ratio
  $v_{\rm bare}/m_e = \varphi^{27}$ is a structural identification
  at the bare rung level.  The gap-function correction
  $\mathrm{gap}(v) \approx 0.26$ accounts for the remaining $9\%$
  to the measured $v = 246.22\;\mathrm{GeV}$.

  \item \textbf{Falsifier}: Measurement of $v/m_e$ inconsistent
  with $\varphi^{27\text{--}28}$ at $> 5\sigma$ would challenge the
  $\varphi$-ladder rung assignment.
\end{enumerate}

%% ============================================================
\section{Conclusion}
\label{sec:conclusion}
%% ============================================================

The electroweak VEV sits at rung 29 of the $\varphi$-ladder:
$v_{\rm bare} = \varphi^{27} m_e \approx 224\;\mathrm{GeV}$ (bare rung,
$9\%$ below the measured $246\;\mathrm{GeV}$; the gap-function
correction closes this residual, as discussed in
Remark~\ref{rem:gap}).  The hierarchy problem is dissolved:
$v/M_{\mathrm{Pl}} = \varphi^{-80} \approx 2.0 \times 10^{-17}$,
an 80-rung gap on the $\varphi$-ladder with no fine-tuning.
From $v$, the $W$, $Z$, and Higgs masses follow with no free
parameters.

%% ============================================================
\begin{thebibliography}{99}

\bibitem{RS_Axioms}
J.~Washburn, ``The Algebra of Reality: A Recognition Science Derivation
of Physical Law,'' \emph{Axioms} \textbf{15}(2), 90 (2026).
\texttt{doi:10.3390/axioms15020090}

\bibitem{Aczel1966}
J.~Acz\'el, \emph{Lectures on Functional Equations and Their
Applications} (Academic Press, 1966).

\bibitem{PDG2024}
S.~Navas \textit{et al.}\ (Particle Data Group),
``Review of Particle Physics,'' Phys.\ Rev.\ D \textbf{110}, 030001
(2024).

\end{thebibliography}

\end{document}
